\scanpage{43}
On trouve ici que
\[
 \beta_\bullet(C^i)=\deleted{\unsure{C^i_\bullet}}
 (\ldots 0\to\ldots 0\to\overset{\deleted{\unsure{i-1|i|i+1}}i}{k}
 \to\overset{\deleted{\unsure{i|i+1}}i+1}{k}\to0\to\ldots\to0)
 \quad\inserted{\text{pour }i\geq1}.
\]
\inserted{$\beta_\bullet(C^0)=(\ldots\deleted{\underset{\text{deg }0}{k}}\to0\to
 \underset{\text{deg }0}{k}\to0)$}
et $\beta_\bullet(d^i): C^i_\bullet\xrightarrow{d^i_\bullet} C^{i+1}_\bullet$
est donné par
\[
 (d^{\deleted{\unsure{i+1}}i}_{\deleted{\unsure{i+1|i}}\bullet})=\mathrm{id}_k
 \qquad \unsure{i\geq10|i\geq0}
\]
\deleted{$d\qquad\unsure{i-1|i=0}$}

% slanted
\margin{Ce n'est pas bien bon à cause de l'irrégularité pour \unsure{ces couples} plutôt}
\deleted{On \unsure{donne|a donc} \unsure{comme|ainsi} \unsure{type|un} modèle la cat.
(\unsure{des} \inserted{$/k$}) complexes de chaînes \unsure{déj. conn.} que nous
\unsure{avions|considéré} comme possibilité de construction d'un couple $(C,C^\bullet)$.}

\uline{Passage de l'exemple 1 à l'exemple 2}

Considérons le foncteur
\[
\begin{aligned}
 M&\longmapsto M^\vee\\
 \mathrm{Ab}_k^\circ&\unsure{\longrightarrow|-}\mathrm{Ab}_k
\end{aligned}
\]
qui induit une équivalence de la catégorie des \inserted{$k$-}modules libres de t.f. avec elle-m.
Cela donne aussi une équivalence entre les cat de complexes de chaînes et cochaînes
à \unsure{comp.} \deleted{\unsure{$k^\vee$ de t.f.}} libres de t.f., et en
celles où les $Z_n$ (et $Z^i$) sont libres de t.f. et \unsure{h. libres|les bords}.
\deleted{\unsure{sont libres, en sont tous}. On \unsure{peut}, \unsure{ce qui du point de vue naturel}}
\deleted{\unsure{pour (pos)}} \unsure{a donc} \unsure{le produit}
\inserted{$\alpha_{\bullet}(\unsure{K|x|\mathbb K})\overset{\unsure{\text{défn}|\text{de}|\text{dual}}}{=}\alpha^\bullet(x)^\vee$} \deleted{\unsure{donné}}
un diagramme d'équivalences

\diagram{\begin{tikzcd}[ampersand replacement=\&]
 \deleted{\unsure{K_\bullet|\mathbb K_\bullet}}\& C\arrow[r,"{\alpha_\bullet}","\sim"']\arrow[d,"{(\alpha^\bullet)^\circ}"',"\sim"] \& K_\bullet\\
 \& K^{\bullet\circ}\arrow[ur,dash,"\sim"] \&
\end{tikzcd}}

% margin note written sideways
\margin{donc $\alpha_n(x)\simeq\mathrm{Hom}(C^{n-1},x)$ si $n\geq1$
et les hom $\alpha_n(x)\to\alpha_{n-1}(x)$ pour $n\geq2$
sont donnés par transposition à partir de $C^{n-1}\to C^n$.
Mais $\alpha_0(x)\simeq\mathrm{Hom}(x,C^0)^\vee$
et $\alpha_1(x)\to\alpha_0(x)$ donné par
$\mathrm{Hom}^{\unsure{\vee}}(C^0,x)\to\mathrm{Hom}(x,C^0)^\vee$ \unsure{ce qu'on aura|qu'on devine}.}

\deleted{Il suffit de définir un isom pour les restrictions à $A$ ce qui est trivial.}

\uline{Explicitons \unsure{quelque peu}.}

\uline{Prop 2} On a une dualité parfaite fonctorielle en $M_\bullet\in\mathrm{Ob}\,K_\bullet$
(\unsure{à|de} $n\geq0$ fixé)
\[
\begin{aligned}
 \alpha_n(\unsure{M|x})_{\unsure{\bullet}}\ \times\ 
 \alpha^n(M)&\longrightarrow k\\[-2pt]
 \rotatebox{90}{$=$}\qquad \rotatebox{90}{$=$} &\qquad \rotatebox{90}{$\simeq$}\\
 \mathrm{Hom}(C^n,x)\times\mathrm{Hom}(x,C^{n+1})&\longrightarrow
 \mathrm{Hom}(C^n,C^{n+1})\deleted{\unsure{\to}}\simeq k\\
 u,v&\longmapsto\unsure{vu}.
\end{aligned}
\]
% archive sheet number: 42
